\chapter{Mérések}

Jelen fejezet a kísérleti mérések végrehajtását és a nyers adatok gyűjtésének folyamatát írja le.
A mérések célja a CRYSTALS-Kyber és a BIKE algoritmusok futási idejének, memóriaigényének
és --- ahol alkalmazható --- energiafogyasztásának meghatározása a korábban bemutatott
tesztkörnyezetben.

%----------------------------------------------------------------------------
\section{Mérési forgatókönyvek}
%----------------------------------------------------------------------------

\subsection{Alap kriptográfiai műveletek mérése}

% TODO: Kulcsgenerálás (KeyGen), kapszulázás (Encapsulation) és visszafejtés (Decapsulation)
% futási idejének izolált mérése — mindkét algoritmuson, különböző biztonsági szinteken
% (Kyber-512/768/1024, BIKE-L1/L3/L5).

\subsection{TLS kézfogás (handshake) mérése}

% TODO: Teljes TLS 1.3 kézfogás időmérése PQC KEM-mel kombinálva (hibrid és tiszta PQC mód).
% MQTT connect–publish–disconnect ciklus mérése.

\subsection{Terheléses mérés}

% TODO: Ismételt kézfogások szekvenciálisan és párhuzamosan — átviteli sebesség és
% átlagos handshake-idő meghatározása terhelés alatt.

%----------------------------------------------------------------------------
\section{Mérési módszer és adatgyűjtés}
%----------------------------------------------------------------------------

\subsection{Időmérés megvalósítása}

% TODO: Mérési könyvtár / syscall (clock_gettime, CLOCK_MONOTONIC), mérési iterációk száma,
% statisztikai módszer (átlag, medián, szórás, kiugróértékek kizárása).

\subsection{Memóriahasználat mérése}

% TODO: Peak RSS / stack usage mérése (valgrind massif, /proc/\$PID/status, vagy
% wolfSSL beépített memóriastatisztika).

\subsection{Energiafogyasztás mérése}

% TODO: Ha elvégezték: mérési pont, mintavételezési frekvencia, mértékegység (mJ/művelet).
% Ha nem: alfejezet elhagyható.

%----------------------------------------------------------------------------
\section{Mérési eredmények}
%----------------------------------------------------------------------------

\subsection{Kyber-512 eredmények}

% TODO: KeyGen / Encap / Decap futási idő táblázat (ms), memória (kB), iterációk száma.

\subsection{Kyber-768 és Kyber-1024 eredmények}

% TODO: Ugyanaz magasabb biztonsági szinteken.

\subsection{BIKE-L1 eredmények}

% TODO: KeyGen / Encap / Decap futási idő, memória. Megjegyzés a dekódolási hibaarányról
% (dekódolási valószínűség, esetleges újrapróbálkozások).

\subsection{BIKE-L3 és BIKE-L5 eredmények}

% TODO: Ugyanaz magasabb szinteken.

\subsection{TLS kézfogás összesítő adatok}

% TODO: Teljes handshake idő algoritmustípusonként, összehasonlítás RSA/ECDH referenciával.

%----------------------------------------------------------------------------
\section{A mérések összefoglalása}
%----------------------------------------------------------------------------

% TODO: Rövid szöveges összegzés a nyers adatokról, kiemelve a legfontosabb tendenciákat
% (melyik algoritmus volt gyorsabb / erőforrás-hatékonyabb), részletes elemzés nélkül
% (az a következő fejezetbe kerül).
